%!TEX program = xelatex
\documentclass[12pt, a4paper, twoside]{scrartcl}

% Die Standardschrift Computer Modern hat keine Umlaute
\usepackage{fontspec}
\setmainfont{Adobe_Caslon_Pro.otf}
%\setmainfont{Theano Didot}

% Polyglossia ist die XeTeX-Alternative zu Babel
\usepackage{polyglossia}
\setdefaultlanguage{german}

% Für den Zeilenabstand
\usepackage{setspace}
\onehalfspacing

% Mathe-Umgebungen
\usepackage{amsmath}

% Absätze werden durch Einrückung der ersten Zeile gekennzeichnet. Die Größe der Einrückung sollte 
% dabei mindestens der Punktgröße der Schrift entsprechen, das vierfache dieser Größe jedoch nicht
% überschreiten. Da wir unser Dokument mit einer Schriftgröße von 12pt setzen, ergibt sich also eine
% Spanne von 12-48pt. Da der Text über die gesamte Breite der A4-Seite gesetzt ist, wählen wir einen
% Wert am oberen Ende der Skala.
%\usepackage{setlength}
\setlength{\parindent}{20pt}

% Mikrotypographische Verfeinerungen des Textsatzes aktivieren.
%\usepackage{microtype}

% Tabularx hat bessere Funktionen für Tabellen als die tabular-Umgebung
\usepackage{tabularx}
\usepackage{array}

% Die Symbole brauchen wir für die Ankreuzkästen bei den Items im Anhang
\usepackage{latexsym} 

% Hier komme die Pakete für die Bibliographie. Als Backend wird biber verwendet. Daher nach dem Compilen
% der TeX-Datei einmal biber auf die .bcf-Datei aufrufen und dann nochmal compieln. Dann sollte alles 
% klappen. 
\usepackage[autostyle]{csquotes}

% Für Tabellen die länger als eine Seite sind
\usepackage{longtable}

\setlength{\headheight}{30pt}

% Kopfzeilen und Fußzeilen
\usepackage{fancyhdr}
\pagestyle{fancy}
\fancyhf{}
\fancyhead[LE,RO]{\nouppercase \leftmark}
\fancyhead[RE,LO]{\rightmark}
\fancyfoot[CE,CO]{}
\fancyfoot[LE,RO]{\thepage}

\usepackage[
    backend=biber,
    bibencoding=utf8,
    %style=numeric-verb,
    %style=authoryear-ibid,
    %style=chicago-authordate,
    %style=apa,
    style=authoryear,
    sortlocale=de_DE,
    maxcitenames=2,
    maxbibnames=99
]{biblatex}

\addbibresource{sample.bib}
\setlength{\bibitemsep}{12pt}
\renewcommand{\labelnamepunct}{:\space}

% Hier wird u.a. durch et al. ersetzt, sowie "und" bei zwei Autoren durch "&"
\DefineBibliographyStrings{german}{
	andothers = {{et\,al\adddot}},
	and = \&,
}

\renewcommand*{\postnotedelim}{\addcolon\space}

% Das entfernt die Anführungszeichen um den Titel in der Bibliographie
\DeclareFieldFormat[inproceedings]{title}{#1} 
\DeclareFieldFormat[incollection]{title}{#1} 
\DeclareFieldFormat[misc]{title}{#1} 
\DeclareFieldFormat[article]{title}{#1} 
\DeclareFieldFormat[book]{title}{#1} 
\DeclareFieldFormat[phdthesis]{title}{#1} 

% Normalerweise schreibt es bei den Namen im Literaturverzeichnis Nachname, Vorname, Vorname, Nachname, Vorname, Nachname.
% Das ist natüröich quatsch
\DeclareNameAlias{sortname}{last-first}

% Das Paket brauchen wir um die Titelseite einzufügen
\usepackage{pdfpages}

\begin{document}

%
% Titelseite in Latex
%
\iffalse
\subject{\vspace{5cm} Seminararbeit}
\title{Masterarbeit Vorlage}
\author{Lieschen Müller}
\date{}
\maketitle
\thispagestyle{empty}
\newpage
\fi


%
% Abstract
%
%\vspace*{\fill}
%\begin{center}
%	{\Large \bfseries \sffamily Abstract} \\
%	\input{Text/Abstract}
%\end{center}
%\vspace*{\fill}
%\thispagestyle{empty}
%\newpage

%
% Inhaltsverzeichnis
%
\tableofcontents
\thispagestyle{empty}
\newpage

%
% Der Text. Hier werden die Dateien aus dem Verzeichnis ./Text eingebunden.
%
\input{Text/Kapitel1}

\newpage

%
% Bibliographie
%
\newpage

% Normalerweise würde das Literaturverzeichnis nicht im Inhaltsverzeichnis auftauchen, das 
% wird mit dieser Zeile behoben.
\addcontentsline{toc}{section}{Literaturverzeichnis}

% Zum Test: Nocite zeigt alle Einträge an, auch wenn sie im Text nicht referenziert werden.
\nocite{*}
\printbibliography
\newpage

\addcontentsline{toc}{section}{Erklärung zur Eigenständigkeit}
\includegraphics[width=\textwidth]{Abbildungen/home-cat.jpg}
\newpage


\addcontentsline{toc}{section}{Anhang}
% Der Anhang
\appendix


\pagestyle{fancy}
\fancyhf{}
\fancyfoot[R]{\thepage}
\fancyhead[L]{\nouppercase{\rightmark}}

\thispagestyle{empty}
\vspace*{\fill}
\begin{center}
    
    {\Huge \sffamily \bfseries Anhang}

\end{center}
\vspace*{\fill}
\newpage

\section{Anhangskapitel}
TestTestttst
%\renewcommand{\thepage}{\arabic{page}}
%\renewcommand{\thepage}{A-\arabic{page}}


\end{document}